% !TeX root = ./../../main.tex
% Techniki konstrukcji obwodów (podrozdział 2.1, rys. fig:tb01:techniques).
% (a) przeniesienie fazy (phase kickback) z interferencją na kubicie kontrolnym,
% (b) dekomputacja (uncomputation) według Bennetta: U_f, kopia wyniku bramką CNOT, U_f^dagger,
% (c) rozkład bramki CP(phi) na bramki CNOT i bramki przesunięcia fazy (sprawdzony numerycznie):
%     CP(phi) = [I x P(phi/2)] CNOT [I x P(-phi/2)] CNOT [P(phi/2) x I].
\begin{subfigure}[b]{0.27\textwidth}
  \centering
  \begin{quantikz}[column sep=0.18cm, row sep={0.7cm,between origins}]
    \lstick{$\ket{0}$} & \gate{\mathrm{H}} & \ctrl{1} & \gate{\mathrm{H}} & \meter{} \\
    \lstick{$\ket{u}$} & & \gate{U} & &
  \end{quantikz}
  \caption{}\label{fig:tb01:techniques:kickback}
\end{subfigure}\hfill
\begin{subfigure}[b]{0.27\textwidth}
  \centering
  \begin{quantikz}[column sep=0.18cm, row sep={0.5cm,between origins}]
    \lstick{$\ket{x}$} & \gate[3]{U_f} & & \gate[3]{U_f^\dagger} & \\
    \lstick{$\ket{0}$} & & \ctrl{2} & & \\
    \lstick{$\ket{0}$} & & & & \\
    \lstick{$\ket{0}$} & & \targ{} & &
  \end{quantikz}
  \caption{}\label{fig:tb01:techniques:uncompute}
\end{subfigure}\hfill
\begin{subfigure}[b]{0.44\textwidth}
  \centering
  \begin{quantikz}[column sep=0.15cm, row sep={0.7cm,between origins}]
    & \ctrl{1} & \\
    & \phase{\varphi} &
  \end{quantikz}
  $=$
  \begin{quantikz}[column sep=0.12cm, row sep={0.7cm,between origins}]
    & \gate{\mathrm{P}(\frac{\varphi}{2})} & \ctrl{1} & & \ctrl{1} & & \\
    & & \targ{} & \gate{\mathrm{P}(-\frac{\varphi}{2})} & \targ{} & \gate{\mathrm{P}(\frac{\varphi}{2})} &
  \end{quantikz}
  \caption{}\label{fig:tb01:techniques:transpile}
\end{subfigure}
